\section{Proof Sketches and Artifact Details}
\label{app:proofs}

\subsection{Proof Dependency Map}
\label{app:proof-map}

The formal statements in Section~\ref{sec:properties} rely on two ingredients: (1)~the local case split induced by the translation rules in Table~\ref{tab:translation-rules}, and (2)~the invariant that the encoding $\mathsf{enc}$ preserves the control-location, resource-count, waiter-count, and variable-store components that are relevant to bug detection and goal checking. The proof obligations therefore decompose into per-operation lemmas plus short inductions over firing sequences.

\paragraph{Forward simulation (Theorem~\ref{thm:forward}).}
For each \textsc{Cir} operation, the translator emits either one \textsc{Cvn} transition or a fixed finite bundle of transitions anchored to the same $\mathit{sid}$. Sequential operations preserve the control token, synchronization operations additionally consume or produce the matching resource-place tokens, and data operations apply the same store update through~$U$. The condvar case is the only non-atomic one: \textsf{wait} decomposes into WaitEnter, Wake1/WakeA, and Reacquire, whose composition preserves the \textsc{Cir} pre/postcondition exactly.

\paragraph{Backward simulation under concrete valuation (Theorem~\ref{thm:backward}).}
When no variable is mapped to~$\top$, every guard evaluates to a concrete boolean. Enabledness in the \textsc{Cvn} is then equivalent to satisfaction of the corresponding \textsc{Cir} precondition, because the input-arc test matches resource availability and the guard test matches branch or data predicates. The only additional state carried by the \textsc{Cvn} consists of condvar auxiliaries, which are updated by literal arithmetic and boolean assignments and therefore remain concrete as long as the original store is concrete.

\paragraph{Deadlock and goal soundness (Theorems~\ref{thm:deadlock-corr} and~\ref{thm:goal-soundness}).}
Deadlock correspondence follows once the enabled sets coincide under the encoding. Goal soundness follows from the same simulation relation plus the observation that each goal predicate is interpreted directly on the encoded marking or valuation: function completion corresponds to the return control place, resource availability corresponds to the initial token count of the matching resource place, and variable equality is checked on the same concrete value carried by both semantics.

\subsection{Artifact and Reproducibility Details}
\label{app:artifact}

The implementation artifact consists of the Rust crate \texttt{cir2cvn}, the benchmark CIRs under \texttt{tests/e2e}, and the experiment runner under \texttt{experiments/}. The core command-line entry points are:
\begin{itemize}[leftmargin=1.4em]
  \item \texttt{cargo build --release}
  \item \texttt{target/release/cir2cvn --validate <file.json>}
  \item \texttt{target/release/cir2cvn --analyze <file.json>}
  \item \texttt{target/release/cir2cvn --goals <file.json>}
  \item \texttt{python experiments/run\_experiment.py --config experiments/config.toml}
\end{itemize}

The configuration file \texttt{experiments/config.toml} fixes the experimental budgets and model parameters used in the paper: $K_{\mathit{gen}}=5$, $K_{\mathit{rep}}=5$, temperature $=0$, output limit $=4096$ tokens, and 3 repeated runs per pattern. It also records the benchmark directories, output directory, and the API endpoints used for the five evaluated models. The local compute environment for the reported experiments is Apple~M4~Pro with 24\,GB RAM. The analysis backend is fully local; model inference is performed through remote APIs configured by environment variables (\texttt{ALL\_API\_KEY}, \texttt{DEEPSEEK\_API\_KEY}, and \texttt{DASHSCOPE\_API\_KEY}).

For the local back-end, the full state-space exploration for every benchmark completes in at most 1\,ms (Section~\ref{subsec:rq3}). A full experimental sweep consists of 5 models, 9 generation tasks, 6 repair tasks, and 3 repetitions per task under the round budgets above. We report exact benchmark outcomes and mean wall-clock times, while provider-side accelerator details and total remote compute consumption remain opaque because the evaluated LLMs are accessed as hosted services rather than self-hosted weights.
